\documentclass[nochapterpage,bigchapter,linedtoc,longdoc,colorback,accentcolor=tud1c]{tudreport}
\usepackage{ngerman}

\usepackage[stable]{footmisc}
\usepackage[ngerman]{hyperref}

\usepackage{longtable}
\usepackage{multirow}
\usepackage{booktabs}

\hypersetup{%
  pdftitle={Entwicklung und Test einer
eingebetteten Elektronik für einen
innovativen Schaltaktor},
  pdfauthor={Malte Breitenbach, Johannes Faupel, Jonas Tautz, Johanna Vetter},
  pdfsubject={Beispieltext},
  pdfview=FitH,
  pdfstartview=FitV
}

%%% Zum Tester der Marginalien %%%
  \newif\ifTUDmargin\TUDmarginfalse
  %%% Wird der Folgende Zeile einkommentiert,
  %%% werden Marginalien gesetzt.
  % \TUDmargintrue
  \ifTUDmargin\makeatletter
    \TUD@setmarginpar{2}
  \makeatother\fi
%%% ENDE: Zum Tester der Marginalien %%%

\newlength{\longtablewidth}
\setlength{\longtablewidth}{0.7\linewidth}
\addtolength{\longtablewidth}{-\marginparsep}
\addtolength{\longtablewidth}{-\marginparwidth}


% \settitlepicture{tudreport-pic}
% \printpicturesize

\title{Entwicklung und Test einer eingebetteten Elektronik für einen innovativen Schaltaktor}
\subtitle{Malte Breitenbach, Johannes Faupel, Jonas Tautz, Johanna Vetter}
%\subsubtitle{email: \textaccent{tud-design@pro-kevin.de}}
%\setinstitutionlogo[width]{TUD_sublogo}
%\uppertitleback{(\textaccent{\textbackslash uppertitleback})}
\lowertitleback{\hfill\today}
\institution{ 
Institut für Mechatronische Systeme im Maschinenbau\\
     Prof. Dr.-Ing. Stephan Rinderknecht}
%\dedication{Hier ist gen"ugend Platz\\
  %f"ur eine Widmung (\textaccent{\textbackslash dedication}).\\
  %\strut\\
  %F"ur Annelore Schmidt\\
  %aus dem Referat Kommunikation.\\
  %Sie hat immer ein offenes Ohr\\
  %f"ur unsere Fragen und Anregungen.}
%\sponsor{\color{tud9b}\rule{\linewidth}{7mm}}
%\sponsor{\hfill\includegraphics[height=6ex]{tud_logo}\hspace{1em}\includegraphics[height=6ex]{TUD_chaos}}

\begin{document}
\maketitle
\begin{abstract}
Hier könnte Ihr Abstract stehen.
\end{abstract}  

\tableofcontents

%% Hier Tex-Dateien includen
%%%%%%%%%%%%%%%%%%%%%%%%%%%%%%%%%%%%%%%%%%
\chapter{Einleitung}
\section{Motivation}


\section{Ziele der Arbeit}
Auch das Institut für Mechatronische Systeme (IMS) der TU Darmstadt nimmt sich dieser Aufgabe an und arbeitet im Rahmen des Projektes Speed4E, dass einen elektrifizierten Antriebsstrang mit Peak-Antriebsdrehzahlen von bis zu $50.000/min$ zum Ziel hat, an einem innovativem Schaltaktor. Im Verlauf vorheriger Arbeiten wurde bereits ein Tauchspulenaktor ausgelegt sowie eine Elektronik für ihn entwickelt. Das Ziel dieser Arbeit ist es nun, die bisherigen Funktionen auf einen Mikrocontroller zu implementieren sowie eine eingebettete Elektronik zu entwerfen, die den Aktor zu einem Smart Actuator transformiert. Des Weiteren sollen Sicherheits- und Überwachungsfunktionen für den Aktor entwickelt werden und Statusmeldungen per CAN gesendet werden. 

\section{Anforderungsliste}
Um das übergeordnete Ziel weiter zu spezifizieren wurde zunächst  eine Anforderungsliste erstellt. In dieser sind alle Forderungen an das Endprodukt gesammelt, sie dient damit als Basis und Referenz für die Produktentwicklung. Die Liste ist hierbei dynamisch, das heißt sie kann im Verlauf des Entwicklungsprozesses verändert oder ergänzt werden.  Die formulierten Anforderungen werden schließlich noch nach Priorität kategorisiert und einer der vier folgenden Anforderungsarten (Quelle: Pahl (2004): Konstruktionslehre - Grundlagen erfolgreicher Produktentwicklung - Methoden und Anwendung S. 189.) zugeteilt:
Festforderungen (FF) sind unter allen Umständen einzuhalten. Eine Erfüllung ist für eine erfolgreiche Lösung notwendig.
Bereichsforderungen (BF) geben einen Toleranzbereich an, innerhalb dessen sich der schlussendlich erreichte Wert befinden muss.
Zielforderungen (ZF) geben an, welcher Wert (auch im Hinsicht auf spätere Entwicklungen) angestrebt wird.
Wünsche (W) sollten nach Möglichkeit erfüllt werden, sind aber keine Voraussetzung. 

\begin{table}[h]
	\centering
		\begin{tabular}{l|p{7cm}|p{7cm}}
			\textbf{Relevanz} & \textbf{Anforderung} & \textbf{Erläuterung} \\ \hline
			& &\\
			FF & Benutzerfreundliche Kommunikation durch CAN Schnittstelle & Empfang von Befehlen, Senden von Statusmeldungen \\
			FF & Nichtflüchtige Kalibrierung & Eine Kalibrierung ist nur einmalig und zur Rekalibrierung notwendig \\
			BF & Schaltzeit & < 100 ms (Latenz zwischen Senden des Befehls und vollständig ausgeführtem Gangwechsel) \\
			FF & Selbstständige Fehlererkennung & Überstrom, Temperatur, Eingangsspannung (OVP/UVP), Dekalibrierung \\
			FF & Schnittstellen & CAN, 8-12VDC Versorgung (max XA), Programmierschnittstelle (für Updates \& Bugfixes) \\
			W & Wartbarkeit & Sicherung wechseln im eingebauten Zustand \\
			BF & kompakte Baugröße & \\ 
			BF & Effizienz (gemittelt über einen Schaltvorgang) & elektrischer Wirkungsgrad > 90 \% \\
			FF & Temperaturbeständigkeit & bis 105°C \\
			BF & Aktorüberschwingen & Toleriert, solange kein unbeabsichtigter Gangwechsel \\
			W & Schaltgabelkraft am Anschlag & möglichst gering \\
			FF & Standby & Standbyleistungsaufnahme < 5W \\
		\end{tabular}
	\caption{Anforderungsliste}
	\label{tab:Anforderungsliste}
\end{table}

\section{Vorgehen der Arbeit}
Nachdem in der oben stehenden Anforderungsliste die Ziele der vorliegenden Arbeit definiert wurden, soll nun das weitere Vorgehen zum Erreichen der Zielsetzung erläutert werden. Zunächst werden in Kapitel 2 der Stand des Prüfstands vor Beginn des ADPs sowie wichtige Grundlagen als Basis für die weitere Bearbeitung und zum besseren Verständnis dargelegt. In Kapitel 3 werden die verschiedenen Bauteile, die für die eingebettete Elektronik benötigt werden, und ihre Funktionen beschrieben. Außerdem wird der erste Prototyp und sein Entwicklungsprozess beschrieben. Kapitel 4 beschäftigt sich mit der Implementierung der Funktionalität in Matlab Simulink BLABLABLA
In Kapitel 5 soll schließlich das endgültige Platinendesign vorgestellt werden.
%%%%%%%%%%%%%%%%%%%%%%%%%%%%%%%%%%%%%%%%%%

\listoffigures\addcontentsline{toc}{chapter}{\listfigurename}
\end{document}
